\documentclass[11pt]{article}
\usepackage[margin=1in]{geometry}
\usepackage{amsmath,amssymb,amsthm}
\usepackage[colorlinks=true,urlcolor=blue,linkcolor=blue]{hyperref}
\newtheorem{theorem}{Theorem}
\newtheorem{lemma}{Lemma}
\newcommand{\D}{\mathbb D}
\newcommand{\Rea}{\operatorname{Re}}
\title{Zero-free sections of square roots of geometric sums\\
\large A solution of MathDB 369171}
\author{Principia Math In Progress}
\date{11 September 2026}
\begin{document}
\maketitle

\section*{Task 1. Prove the zero-free conjecture}

For an integer $q\geq1$, use the branch equal to $1$ at the origin and set
\[
 f_q(z)=\sqrt{\frac{1-z^q}{1-z}}
       =\sum_{k\geq0}b_kz^k,\qquad
 S_n(z)=\sum_{k=0}^n b_kz^k.
\]
Kovalev's Conjecture 1.2~\cite{K} is the assertion below for $q=d+1\geq4$.
We prove a stronger assertion, including strict control on the boundary.

\begin{theorem}\label{thm:main}
For every integer $q\geq1$ and $n\geq0$,
\[
 \Rea\bigl(\sqrt{1-z}\,S_n(z)\bigr)>0
 \quad (|z|\leq1,\ z\ne1),
 \qquad S_n(1)>0.
\]
Consequently $S_n$ has no zero in the closed unit disk.
\end{theorem}

All square roots below are principal ones. In particular they are
unambiguous for $1-z$ and $1-tz$, $|z|\leq1$, $0\leq t<1$.
The proof uses a remainder represented by a signed beta density whose
negative and positive parts have ordered supports.

\subsection*{Coefficients and a normalized remainder}

Put
\[
 c_k=\frac{\binom{2k}{k}}{4^k},\qquad
 a_j=\frac{c_j}{2j-1}\quad(j\geq1),\qquad
 A_m=(2m+1)c_m,\qquad C_m(z)=\sum_{k=0}^m c_kz^k.
\]
The binomial expansion and telescoping give
\begin{align}
 \sqrt{1-w}&=1-\sum_{j\geq1}a_jw^j, \notag\\
 \sum_{k=0}^m c_k&=A_m,\qquad
 \sum_{j=1}^M a_j=1-c_M,\qquad
 \sum_{j=1}^M ja_j=M c_M. \label{eq:sums}
\end{align}
For the last identity, $ja_j=c_{j-1}/2$ and
$A_{M-1}=2Mc_M$ when $M\geq1$.
Also $A_{m+1}/A_m=(m+3/2)/(m+1)>1$.

Let $M=\lfloor n/q\rfloor$. Coefficient convolution yields
\begin{equation}
 S_n(z)=C_n(z)-\sum_{j=1}^M a_j z^{qj} C_{n-qj}(z).
 \label{eq:section}
\end{equation}
At $z=1$ this already proves
\begin{equation}
 S_n(1)=A_n-\sum_{j=1}^M a_j A_{n-qj}
       \geq A_n c_M>0. \label{eq:one}
\end{equation}

Define $I_m(z)=1-\sqrt{1-z}\,C_m(z)$. Differentiating, and using
$(k+1)c_{k+1}=(k+1/2)c_k$, gives
\[
 I_m'(z)=\frac{A_m}{2}\frac{z^m}{\sqrt{1-z}},\qquad I_m(0)=0.
\]
It follows that
\begin{equation}
 I_m(z)=z^{m+1}\int_0^1 p_m(t)u_z(t)\,dt,\qquad
 p_m(t)=\frac{A_m}{2}\frac{t^m}{\sqrt{1-t}},\qquad
 u_z(t)=\sqrt{\frac{1-t}{1-tz}}.
 \label{eq:beta}
\end{equation}
Each $p_m$ is a probability density on $(0,1)$: integration of
$I_m'$ from $0$ to $1$ gives $\int_0^1p_m=I_m(1)=1$.
The integral identity extends continuously to $|z|\leq1$.
At $z=1$, set $u_1(t)=1$ for $t<1$; its value at the endpoint is immaterial.

\subsection*{An ordered-kernel inequality}

\begin{lemma}\label{lem:kernel}
For $|z|\leq1$ and $0\leq s\leq t<1$,
\[
 |u_z(t)|\leq1,\qquad |2u_z(t)-u_z(s)|\leq |u_z(s)|\leq1.
\]
If $z\ne1$ and $t>0$, then $|u_z(t)|<1$.
\end{lemma}
\begin{proof}
Put $w=1-z$, $a=s/(1-s)$ and $b=t/(1-t)$.
Then $\Rea w\geq0$, $0\leq a\leq b$, and
\[
 u_z(s)=(1+aw)^{-1/2},\qquad
 u_z(t)=(1+bw)^{-1/2}.
\]
The modulus bounds follow from $|1+aw|\geq1$.
For the remaining inequality set $r=(1+aw)/(1+bw)$. Direct calculation gives
\[
 \Rea r-|r|^2
 =\frac{(b-a)(\Rea w+a|w|^2)}{|1+bw|^2}\geq0.
\]
In particular $|r|\leq1$. For $v=\sqrt r=u_z(t)/u_z(s)$, the branch
is the principal one, since both denominator arguments lie in
$[-\pi/2,\pi/2]$ and their difference lies in that interval. Thus
\[
 (\Rea v)^2=\frac{|r|+\Rea r}{2}
       \geq\frac{|r|+|r|^2}{2}\geq |r|^2.
\]
So $\Rea v\geq |r|=|v|^2$, whence $|2v-1|\leq1$.
Multiplication by $|u_z(s)|$ proves the claim.
For $|z|\leq1$, $z\ne1$, one has $\Rea w>0$, making
$|1+bw|>1$ when $t>0$.
\end{proof}

\subsection*{The signed density has small negative mass}

For $M\geq1$, define
\[
 h(t)=p_n(t)-\sum_{j=1}^M a_j p_{n-qj}(t),
 \qquad c=c_M.
\]
Its integral is $c$. Write $\lambda_j=A_{n-qj}/A_n$.
The ratio
\begin{equation}
 \frac{h(t)}{p_n(t)}
     =1-\sum_{j=1}^M a_j\lambda_j t^{-qj}
 \label{eq:cross}
\end{equation}
is strictly increasing on $(0,1)$, tends to $-\infty$ at $0$,
and has positive limit $S_n(1)/A_n$ at $1$, by~\eqref{eq:one}.
It therefore has a unique zero $t_0\in(0,1)$.
The density $h$ is negative below $t_0$ and positive above it.

Let $N=\int_0^{t_0}(-h(t))\,dt$ be its negative mass. We claim
\begin{equation}
 0<N<c. \label{eq:mass}
\end{equation}
Indeed, $p_{n-qj}(t)\geq\lambda_j p_n(t)$ for $0<t<1$.
Consequently
\[
 h=\left(1-\sum_j a_j\lambda_j\right)p_n
          -\sum_j a_j\bigl(p_{n-qj}-\lambda_jp_n\bigr),
\]
where the coefficient of $p_n$ and all densities in parentheses are
nonnegative. The negative part of $h$ is therefore bounded by the last sum,
so
\[
 N\leq \sum_{j=1}^M a_j(1-\lambda_j).
\]
For $0\leq m<n$,
\[
 \frac{A_m}{A_n}
 =\prod_{k=m+1}^n\frac{k}{k+1/2}
 >\prod_{k=m+1}^n\frac{k-1}{k}=\frac mn.
\]
Taking $m=n-qj$ and using $n\geq qM$ gives
$\lambda_j>1-qj/n\geq1-j/M$. Therefore~\eqref{eq:sums} implies
\[
 N<\frac1M\sum_{j=1}^M ja_j=c_M=c,
\]
as claimed.

\subsection*{Remainder bound and completion of the proof}

Write
\[
 G_M(z)=1-\sum_{j=1}^M a_jz^{qj},\qquad
 D_n(z)=I_n(z)-\sum_{j=1}^M a_jz^{qj}I_{n-qj}(z).
\]
Equations~\eqref{eq:section} and~\eqref{eq:beta} give
\begin{equation}
 \sqrt{1-z}\,S_n(z)=G_M(z)-D_n(z),\qquad
 D_n(z)=z^{n+1}\int_0^1 h(t)u_z(t)\,dt. \label{eq:remainder}
\end{equation}
We show that
\begin{equation}
 |D_n(z)|<c_M\qquad (|z|\leq1,\ z\ne1). \label{eq:strict}
\end{equation}
When $M=0$, take $h=p_n$, $c_M=1$. Then~\eqref{eq:strict} follows
directly from~\eqref{eq:beta} and $|u_z(t)|<1$ for $0<t<1$.

For $M\geq1$, normalize the positive and negative parts of $h$ to
probability measures $P$ and $Q$, respectively. Their supports are contained
in $(t_0,1)$ and $(0,t_0)$, and their masses before normalization are
$c+N$ and $N$. Set $B=N/c$, so $0<B<1$ by~\eqref{eq:mass}.
Using independent variables $t$ with law $P$ and $s$ with law $Q$,
\begin{align*}
 \frac1c\int h(t)u_z(t)\,dt
  &=(1+B)\,\mathbb E_Pu_z(t)-B\,\mathbb E_Qu_z(s)\\
  &=(1-B)\,\mathbb E_Pu_z(t)
       +B\,\mathbb E_{P\otimes Q}\bigl(2u_z(t)-u_z(s)\bigr).
\end{align*}
Here $s<t$ almost surely. Lemma~\ref{lem:kernel} shows that the second
expectation has modulus at most $1$. The first has modulus strictly
less than $1$ when $z\ne1$, since $t>t_0>0$.
As $1-B>0$, the displayed expression has modulus strictly less than $1$.
Multiplying by $c|z|^{n+1}\leq c$ proves~\eqref{eq:strict}.

Finally, \eqref{eq:sums} implies
\[
 \Rea G_M(z)\geq1-\sum_{j=1}^M a_j|z|^{qj}\geq c_M
 \qquad (|z|\leq1).
\]
Combining this with~\eqref{eq:remainder} and~\eqref{eq:strict} proves
$\Rea(\sqrt{1-z}S_n(z))>0$ for $z\ne1$. At $z=1$,
\eqref{eq:one} applies. Since $\sqrt{1-z}\ne0$ away from $z=1$,
the asserted zero-freeness follows. This proves Theorem~\ref{thm:main}.
\qed

\paragraph{Scope.}
This is an all-parameter analytic proof; finite computations are supplementary.
The distinct assertion $\Rea S_n(e^{it})>0$ is not needed here.
That assertion is Conjecture 1.3 of~\cite{K} (MathDB 369172), for which a
separate exact counterexample has been saved.

\section*{Task 2. Determine sharp segment bounds}
Use the established zero-freeness in the Landau--Sz\'asz method to determine
the sharp norm of each coefficient segment of fixed length, and classify
the extremal bounded analytic functions.

\section*{Task 3. Quantify the distance of zeros from the disk}
Turn the strict remainder estimate into explicit lower bounds for
$|S_n(z)|$ on the unit circle and explicit lower bounds for the modulus
of its nearest zero, with dependence on both $q$ and $n$.

\section*{Task 4. Classify fractional powers}
Determine the exponents $\alpha>0$ for which every Taylor section of
$((1-z^q)/(1-z))^\alpha$ is zero-free in the closed unit disk for every $q$.
Identify the range in which an ordered signed-density argument still works.

\section*{Task 5. Describe the nearest-zero asymptotics}
Derive uniform asymptotic expansions for zeros approaching the nontrivial
$q$th roots of unity as $n$ grows, and determine how the expansions change
when $q$ grows with $n$.

\begin{thebibliography}{9}
\bibitem{K}
L. V. Kovalev,
\emph{Sharp bounds for some segments of bounded power series},
arXiv:2507.04544v1 (2025), Conjecture 1.2, p.~2.
\url{https://arxiv.org/abs/2507.04544}.
\end{thebibliography}
\end{document}
